\kbvsechead{ವೃತ್ತ ಜಾತಿಯ ಛಂದಸ್ಸಿನ ಲಕ್ಷಣ :}

\vspace{4pt plus 4pt minus 1pt}\noindent{\kbvsecfont\fontsize{11pt}{14.5pt}\selectfont\color{kbvaccentdark} (೧)\hspace{0.45em}ಉತ್ಪಲಮಾಲಾ ವೃತ್ತ :}\hspace{0.5em}ಇದು ನಾಲ್ಕು ಪಾದಗಳ ವೃತ್ತ, ಪ್ರತಿಯೊಂದು ಪಾದದಲ್ಲಿ ಆರು ಗಣಗಳು, ಒಂದು ಲಘು ಮತ್ತು ಒಂದು ಗುರು ಸೇರಿ ಒಟ್ಟು ಇಪ್ಪತ್ತು ಅಕ್ಷರಗಳಿರುವುವು.\par

\begin{kbvnote}
\noindent{\kbvsecfont ವಿನ್ಯಾಸ :}\hspace{0.5em}(ಭ--ರ..-ನ..-ಭ..-ಭ..-ರ..-೧ಲ--೧ಗು)
\end{kbvnote}

\vspace{4pt plus 4pt minus 1pt}\noindent{\kbvsecfont\fontsize{11pt}{14.5pt}\selectfont\color{kbvaccentdark} (೨)\hspace{0.45em}ಚಂಪಕಮಾಲಾ ವೃತ್ತ :}\hspace{0.5em}ಇದು ನಾಲ್ಕು ಪಾದಗಳ ವೃತ್ತ, ಪ್ರತಿಯೊಂದು ಪಾದದಲ್ಲಿ ಏಳು ಗಣಗಳಿದ್ದು ಒಟ್ಟು ಇಪ್ಪತ್ತೊಂದು ಅಕ್ಷರ ಗಳಿರುವುವು.\par

\begin{kbvnote}
\noindent{\kbvsecfont ವಿನ್ಯಾಸ :}\hspace{0.5em}(ನ..-ಜ..-ಭ..-ಜ..-ಜ..-ಜ..-ರ..)
\end{kbvnote}

\vspace{4pt plus 4pt minus 1pt}\noindent{\kbvsecfont\fontsize{11pt}{14.5pt}\selectfont\color{kbvaccentdark} (೩)\hspace{0.45em}ಶಾರ್ದೂಲ ವಿಕ್ರೀಡಿತ ವೃತ್ತ :}\hspace{0.5em}ಇದು ನಾಲ್ಕು ಪಾದಗಳ ವೃತ್ತ. ಪ್ರತಿಯೊಂದು ಪಾದದಲ್ಲಿ ಆರು ಗಣಗಳು ಮತ್ತು ಒಂದು ಗುರು ಸೇರಿ ಒಟ್ಟು ಹತ್ತೊಂಬತ್ತು ಅಕ್ಷರಗಳಿರುವುವು.\par

\begin{kbvnote}
\noindent{\kbvsecfont ವಿನ್ಯಾಸ :}\hspace{0.5em}(ಮ--ಸ..-ಜ..-ಸ..-ತ..-ತ..-೧ಗು)
\end{kbvnote}

\vspace{4pt plus 4pt minus 1pt}\noindent{\kbvsecfont\fontsize{11pt}{14.5pt}\selectfont\color{kbvaccentdark} (೪)\hspace{0.45em}ಮತ್ತೇಭ ವಿಕ್ರೀಡಿತ ವೃತ್ತ :}\hspace{0.5em}ಇದು ನಾಲ್ಕು ಪಾದಗಳ ವೃತ್ತ. ಪ್ರತಿಯೊಂದು ಪಾದದಲ್ಲಿ ಆರು ಗಣಗಳು, ಒಂದು ಲಘು ಮತ್ತು ಒಂದು ಗುರು ಸೇರಿ ಒಟ್ಟು ಇಪ್ಪತ್ತು ಅಕ್ಷರಗಳಿರುವುವು.\par

\begin{kbvnote}
\noindent{\kbvsecfont ವಿನ್ಯಾಸ :}\hspace{0.5em}(ಸ..-ಭ..-ರ..-ನ..-ಮ..-ಯ..-೧ಲ--೧ಗು)
\end{kbvnote}

\vspace{4pt plus 4pt minus 1pt}\noindent{\kbvsecfont\fontsize{11pt}{14.5pt}\selectfont\color{kbvaccentdark} (೫)\hspace{0.45em}ಸ್ರಗ್ಧರಾ ವೃತ್ತ :}\hspace{0.5em}ಇದು ನಾಲ್ಕು ಪಾದಗಳ ವೃತ್ತ. ಪ್ರತಿ ಯೊಂದು ಪಾದದಲ್ಲಿ ಏಳು ಗಣಗಳಿದ್ದು ಒಟ್ಟು ಇಪ್ಪತ್ತೊಂದು ಅಕ್ಷರಗಳಿರು ವುವು.\par

\begin{kbvnote}
\noindent{\kbvsecfont ವಿನ್ಯಾಸ :}\hspace{0.5em}(ಮ..-ರ..-ಭ..-ನ..-ಯ..-ಯ..-ಯ..)
\end{kbvnote}
